\noindent\null\hfill{\small\textbf{ZUR HANNOVERSCHEN TRIANGULATION.}}\hfill\makebox[0pt][r]{\fontsize{13}{16}\selectfont \textbf{375}}\par\medskip

\begin{center}
\textsc{Gauss} an \textsc{Olbers}.\quad Göttingen, 25.~Februar 1825.
\end{center}

\par\noindent\hspace*{\parindent}\mbox{.\ .\ .\ .\ .\ .}\quad Jetzt noch ein paar Worte über das geodätische Nivellement.
Ich bin mir über diesen Gegenstand zwar \textls[120]{ganz im Klaren}, weiß aber nicht, ob ich mich in der Kürze so erklären kann, dass Sie mir sogleich zustimmen.

\par\noindent\hspace*{\parindent}Ich habe immer geglaubt, dass der Ausdruck »Localattraction« sehr schlecht gewählt ist und leicht falsche Vorstellungen hervorrufen kann.
Man sollte sagen, dass die Richtungen der Schwere nicht so verlaufen, wie sie bei einem gleichförmigen Sphäroid verlaufen würden.
Die Richtung der Schwere ist die Gesamtwirkung der Anziehung aller Bestandteile des Erdkörpers (und der Zentrifugalkraft), und bei seinem in Bezug auf die Dichte unregelmäßigen Aufbau und den Unebenheiten seiner Oberfläche kann sie nicht dieselbe sein wie bei einem regelmäßigen Sphäroid.
Aber \textls[120]{wie auch der Aufbau sein mag}, immer lässt sich durch jeden Punkt eine Fläche legen, die ganz um die Erde herumreicht und auf der die Richtung der Schwere \textls[120]{genau} senkrecht steht. Die Oberfläche einer zusammenhängenden ruhenden Flüssigkeit würde diese Fläche darstellen.
\textbf{Diese} Fläche ist es, die eine Horizontalfläche heißt (couche de niveau). Den Punkten \textls[120]{dieser} Fläche \textls[120]{weist man dieselbe Höhe zu}, ohne sich im Geringsten darum zu kümmern, ob oder wie weit sie von einem elliptischen Sphäroid abweichen. Die Höhen über dieser Fläche ergeben sich sowohl aus dem Barometer als auch aus der trigonometrischen Messung, sodass beide immer miteinander übereinstimmen müssen.
Dabei wird nur vorausgesetzt\textsuperscript{*}), dass sich auf jeder Dreiecksseite die Richtung der Schwere nach dem Gesetz der Stetigkeit ändert (wenn auch vielleicht schneller oder langsamer als beim elliptischen Sphäroid), und diese Voraussetzung kann nur dann einen kleinen Fehler verursachen, wenn an einer der Dreiecksstationen eine \textls[120]{echte} Lokalattraktion vorliegt, die rein örtlich und auf einen kleinen Raum beschränkt

\par\smallskip\noindent\rule{.24\textwidth}{.4pt}\par\smallskip
{\footnotesize
\noindent\hspace*{.06\linewidth}\textsuperscript{*}) Und natürlich wird auch vorausgesetzt, dass alle Zenitdistanzen reziprok gemessen werden, was bei meiner Messung \textls[120]{ohne Ausnahme} der Fall ist.
Bei \textls[120]{einseitigen} Messungen ist Ihre Bemerkung vollkommen berechtigt, da man dabei die Amplitude sphäroidisch berechnen muss.\par
}

\clearpage
\noindent\makebox[0pt][l]{\fontsize{13}{16}\selectfont \textbf{376}}\hfill{\small\textbf{BRIEFWECHSEL MIT OLBERS.}}\hfill\null\par\medskip
\noindent (außerhalb davon nicht wahrnehmbar) ist.
Ich bin aber überzeugt, dass, höchstens mit Ausnahme des Brockens, im gesamten Bereich meiner Dreiecke und der Dreiecke \textsc{Schumachers} keine \textls[120]{solche} Lokalattraktion vorkommt.\ \mbox{.\ .\ .\ .\ .\ .}
